\chapter{Design and Architecture}

\section*{Introduction}

This chapter is the technical heart of the report. It describes the \caimp{}
microservices architecture, data model, main processing flows and design decisions
that guided development. Every architectural choice is justified against the
alternatives evaluated.

\section{Architectural Choice: Microservices vs. Monolith}

The microservices architecture was selected for three decisive reasons:
(1)~Holt-Winters forecasting and LLM inference are compute-intensive tasks that
must not block the query APIs;
(2)~the Go telemetry writer benefits from goroutines for concurrent ingestion
without Python GIL limitations;
(3)~future metric sources (Kubernetes, AWS CloudWatch) can be added as independent
services without modifying existing ones.

\begin{table}[H]
\centering
\caption{Monolith vs.\ microservices comparison for \caimp{}}
\label{tab:arch-choice}
\small
\begin{tabularx}{\textwidth}{lXX}
\toprule
\textbf{Criterion} & \textbf{Modular Monolith} & \textbf{Microservices} \\
\midrule
Deployment       & Simple (one process)        & Complex (Docker Compose) \\
Scalability      & Vertical only               & Horizontal per service \\
Fault isolation  & One crash affects all       & Graceful degradation \\
Mixed languages  & Impossible                  & Go + Python \\
AI decoupling    & LLM in API process          & Dedicated worker \\
\bottomrule
\end{tabularx}
\end{table}

\section{Global Architecture}

The platform comprises 19 containerised services organised into four functional
layers:
\begin{description}[leftmargin=2cm]
  \item[Ingestion layer] Receives metrics from agents (\servicename{telemetry-writer},
    \servicename{splunk-hec-bridge}) and changes (\servicename{change-tracker}).
  \item[Event layer] NATS JetStream bus providing asynchronous communication.
  \item[Processing layer] Consumer services processing events independently.
  \item[Presentation layer] REST APIs, WebSocket Gateway and React frontend.
\end{description}

\begin{landscape}
\begin{figure}[H]
\centering
\resizebox{\linewidth}{!}{%
\begin{tikzpicture}[
  svc/.style   = {draw=NavyBlue,    fill=NavyBlue!10,     rounded corners=4pt,
                  minimum width=3.2cm, minimum height=0.9cm,
                  font=\normalsize\bfseries, align=center},
  agent/.style = {draw=SuccessGreen, fill=SuccessGreen!12, rounded corners=4pt,
                  minimum width=3.2cm, minimum height=0.9cm,
                  font=\normalsize\bfseries, align=center},
  store/.style = {draw=TealBlue,    fill=TealBlue!10,     rounded corners=4pt,
                  minimum width=3.4cm, minimum height=0.9cm,
                  font=\normalsize\bfseries, align=center},
  bus/.style   = {draw=AccentOrange, fill=AccentOrange!12, rounded corners=6pt,
                  minimum width=15.0cm, minimum height=1.0cm,
                  font=\large\bfseries, align=center,
                  text=AccentOrange!80!black},
  proxy/.style = {draw=DarkGray!70, fill=DarkGray!10,     rounded corners=4pt,
                  minimum width=15.0cm, minimum height=0.9cm,
                  font=\large\bfseries, align=center},
  layer/.style = {font=\small\bfseries\color{NavyBlue!70},
                  anchor=west, inner sep=0},
  arr/.style   = {-{Stealth[length=6pt]}, thick, draw=NavyBlue!60},
  darr/.style  = {-{Stealth[length=6pt]}, thick, draw=NavyBlue!35, dashed},
]
\begin{scope}[on background layer]
  \fill[SuccessGreen!6,  rounded corners=5pt] (-2.4, 8.9) rectangle (13.8, 10.1);
  \fill[NavyBlue!5,      rounded corners=5pt] (-2.4, 7.4) rectangle (13.8,  8.8);
  \fill[AccentOrange!6,  rounded corners=5pt] (-2.4, 5.9) rectangle (13.8,  7.3);
  \fill[NavyBlue!5,      rounded corners=5pt] (-2.4, 4.4) rectangle (13.8,  5.8);
  \fill[TealBlue!5,      rounded corners=5pt] (-2.4, 2.9) rectangle (13.8,  4.3);
  \fill[NavyBlue!5,      rounded corners=5pt] (-2.4, 1.4) rectangle (13.8,  2.8);
  \fill[DarkGray!6,      rounded corners=5pt] (-2.4, 0.1) rectangle (13.8,  1.3);
\end{scope}
% Layer labels
\node[layer] at (-2.3, 9.5)  {Servers};
\node[layer] at (-2.3, 8.1)  {Ingestion};
\node[layer] at (-2.3, 6.6)  {Bus};
\node[layer] at (-2.3, 5.1)  {Processing};
\node[layer] at (-2.3, 3.6)  {Data};
\node[layer] at (-2.3, 2.1)  {Presentation};
\node[layer] at (-2.3, 0.7)  {Access};
% Row 1 — Monitored servers
\node[agent] (a1) at (1.6,  9.5) {Server A};
\node[agent] (a2) at (5.6,  9.5) {Server B};
\node[agent] (a3) at (9.6,  9.5) {Server C};
% Row 2 — Ingestion services
\node[svc] (tw)  at (1.6,  8.1) {Telemetry Writer};
\node[svc] (hec) at (5.6,  8.1) {HEC Bridge};
\node[svc] (ct)  at (9.6,  8.1) {Change Tracker};
% Row 3 — NATS event bus
\node[bus] (nats) at (5.7, 6.6) {NATS JetStream};
% Row 4 — Worker services
\node[svc] (aiw) at (0.8,  5.1) {AI Worker};
\node[svc] (ale) at (3.8,  5.1) {Alert Engine};
\node[svc] (fce) at (7.0,  5.1) {Forecast Engine};
\node[svc] (cor) at (10.5, 5.1) {Correlation Engine};
% Row 5 — Data stores
\node[store] (pg)    at (2.4,  3.6) {PostgreSQL};
\node[store] (redis) at (6.2,  3.6) {Redis :6379};
\node[store] (oll)   at (10.5, 3.6) {Ollama :11434};
% Row 6 — APIs and Frontend
\node[svc] (fe)   at (0.8,  2.1) {Frontend :3000};
\node[svc] (qapi) at (3.8,  2.1) {Query API :8002};
\node[svc] (adm)  at (7.0,  2.1) {Admin API :8001};
\node[svc] (wsg)  at (10.5, 2.1) {WS Gateway};
% Row 7 — Reverse proxy
\node[proxy] (nginx) at (5.7, 0.7) {Nginx --- Reverse Proxy :80 / :443};
% Arrows — Servers → Ingestion
\draw[arr] (a1) -- (tw);
\draw[arr] (a2) -- (hec);
\draw[arr] (a3) -- (ct);
% Arrows — Ingestion → NATS
\draw[arr] (tw)  -- (nats);
\draw[arr] (hec) -- (nats);
\draw[arr] (ct)  -- (nats);
% Arrows — NATS → Workers
\draw[arr] (nats.south) -- (aiw.north);
\draw[arr] (nats.south) -- (ale.north);
\draw[arr] (nats.south) -- (fce.north);
\draw[arr] (nats.south) -- (cor.north);
% Arrows — Workers → Data
\draw[arr] (aiw.south) -- (pg.north);
\draw[arr] (ale.south) -- (redis.north);
\draw[arr] (fce.south) -- (pg.north);
\draw[arr] (cor.south) -- (oll.north);
\draw[darr] (aiw.south) to[out=-70, in=110] (oll.north);
% Arrows — Data → Presentation
\draw[arr]  (pg.south) -- (qapi.north);
\draw[arr]  (pg.south) to[out=-110, in=90] (fe.north);
\draw[arr]  (pg.south) to[out=-80,  in=90] (adm.north);
\draw[darr] (nats.south) to[out=-5, in=90] (wsg.north);
% Arrows — Presentation → Nginx
\draw[arr] (fe.south)   -- (nginx.north);
\draw[arr] (qapi.south) -- (nginx.north);
\draw[arr] (adm.south)  -- (nginx.north);
\draw[arr] (wsg.south)  -- (nginx.north);
\end{tikzpicture}
}%
\caption{Global architecture of the \caimp{} platform --- 7 layers, 19 services}
\label{fig:arch}
\end{figure}
\end{landscape}

\section{How the Architecture Works — End-to-End Flows}

Figure~\ref{fig:arch} shows the static structure of the platform. This section
describes its dynamic behaviour: how data travels through each layer, from
collection on the monitored server to a natural-language explanation appearing
in the operator's browser.

% ─────────────────────────────────────────────────────────────────────────────
\subsection{Flow 1 — Metric Collection and Ingestion}

The \caimp{} agent is a Python daemon deployed on each target server. Every
15 seconds it reads system indicators (CPU, memory, disk) via \texttt{psutil},
encodes them as \textsc{otlp}/HTTP (protobuf), and sends them to the
\servicename{Telemetry Writer} (port 4318). This Go service:

\begin{enumerate}
  \item \textbf{Receives} the OTLP batch and extracts metric data points;
  \item \textbf{Persists} the values to the \texttt{metrics} hypertable
        (TimescaleDB) via a bulk \texttt{COPY} insert;
  \item \textbf{Detects} anomalies in parallel using three algorithms
        (static threshold, Z-score, rate of change);
  \item \textbf{Publishes} each anomaly on the NATS subject
        \texttt{anomaly.detected.\{org\_id\}.\{severity\}}.
\end{enumerate}

In parallel, the agent runs four \textit{change reporters} that continuously
watch Docker events, package installations, SSH sessions and cron jobs. These
events are sent to the \servicename{Change Tracker} (port 8011), which stores
them in \texttt{change\_events} and publishes them on \texttt{change.detected.>}.

% ─────────────────────────────────────────────────────────────────────────────
\subsection{Flow 2 — AI Explanation and Change Correlation}

As soon as an \texttt{anomaly.detected.*} message arrives on the NATS bus,
\textbf{two independent consumer services} wake up:

\paragraph{AI Worker} For each anomaly it:
\begin{enumerate}
  \item Queries PostgreSQL for the last 30 minutes of metric history;
  \item Queries pgvector by cosine similarity to retrieve the most similar
        runbooks and past incidents (\textit{Retrieval-Augmented Generation, \rag{}});
  \item Builds a structured prompt targeting non-DevOps developers and sends
        it to Ollama (Llama 3.1 8B \textsc{llm});
  \item Validates the JSON response, inserts it into \texttt{ai\_explanations},
        then publishes \texttt{ai.explanation.ready.*} on NATS.
\end{enumerate}

\paragraph{Correlation Engine} In parallel it performs causal analysis:
\begin{enumerate}
  \item Queries \texttt{change\_events} in the window $[t-30\text{min},\,t]$
        for the same server;
  \item Scores each candidate change:
        $s = e^{-0.1 \cdot \Delta t} \times w_{\text{type}}$
        (temporal decay $\times$ change-type weight);
  \item Sends the top 3 candidates to Ollama for a causal narrative;
  \item Persists the result in \texttt{root\_cause\_analyses} and publishes
        \texttt{rootcause.ready.*} on NATS.
\end{enumerate}

% ─────────────────────────────────────────────────────────────────────────────
\subsection{Flow 3 — Resource Forecasting}

The \servicename{Forecast Engine} runs autonomously on two periodic loops:

\begin{description}[leftmargin=2cm]
  \item[Forecast loop (every 30 min)] For each server/metric pair, fetches the
    last 6 hours of data, applies Holt-Winters double exponential smoothing,
    generates 24 hourly forecast points with $\pm 1.28\sigma$ confidence
    intervals, computes \textit{time-to-critical}, requests an LLM narrative
    from Ollama, and persists everything in \texttt{server\_forecasts}.
  \item[Evaluation loop (every 60 min)] Compares past forecasts against actual
    values, stores the mean absolute error (\textsc{mae}), and feeds historical
    accuracy into the next LLM prompt to calibrate expressed confidence.
\end{description}

% ─────────────────────────────────────────────────────────────────────────────
\subsection{Flow 4 — Real-Time Browser Updates}

The \servicename{WS Gateway} maintains persistent WebSocket connections with
all connected browsers. It subscribes to:
\texttt{anomaly.detected.>}, \texttt{ai.explanation.ready.>},
\texttt{rootcause.ready.>}, and \texttt{forecast.critical.>}.

When a message arrives, the gateway fan-outs it to all WebSocket clients of
the relevant organisation (multi-tenant isolation enforced). The dashboard
refreshes with no polling — detection-to-display latency is under 3 seconds.

% ─────────────────────────────────────────────────────────────────────────────
\subsection{Flow 5 — Dashboard Query}

When an operator opens the dashboard, the browser sends
\texttt{GET /api/v1/dashboard/summary} to the \servicename{Query API} (via Nginx).
It aggregates server statuses, 24-hour anomalies, at-risk forecasts, and the
latest root-cause analyses into a structured response:
\textit{overall\_status}, \textit{potential\_issues} (sorted by criticality),
and \textit{recommended\_actions} (numbered). The answers-first interface
presents this immediately without requiring the operator to read graphs.

\section{Event Bus: NATS JetStream}

NATS JetStream \cite{nats_docs} was selected over Kafka and RabbitMQ for
its 30\,MB memory footprint, sub-millisecond latency, and simple configuration
while offering the durability guarantees needed (disk persistence, acknowledgements).

\begin{table}[H]
\centering
\caption{NATS JetStream streams in \caimp{}}
\label{tab:nats-streams}
\small
\begin{tabularx}{\textwidth}{llXll}
\toprule
\textbf{Stream} & \textbf{Subjects} & \textbf{Description} & \textbf{Retention} & \textbf{Consumers} \\
\midrule
METRICS     & \texttt{metric.received.>}       & Raw metric batches          & 24h  & 2 \\
ANOMALIES   & \texttt{anomaly.detected.>}      & Detected anomalies          & 7d   & 3 \\
AI\_OUT     & \texttt{ai.explanation.ready.>}  & AI explanations             & 30d  & 1 \\
CHAT        & \texttt{chat.requested.>}        & Chat requests               & 1h   & 1 \\
AUDIT       & \texttt{audit.log.>}             & Audit events                & 90d  & 0 \\
FORECASTS   & \texttt{forecast.>}              & Forecast results            & 30d  & 1 \\
CHANGES     & \texttt{change.detected.>}       & System changes              & 7d   & 2 \\
ROOT\_CAUSE & \texttt{rootcause.ready.>}       & Root cause analyses         & 30d  & 1 \\
\bottomrule
\end{tabularx}
\end{table}

\section{Data Model}

\subsection{Design Principles}

Two fundamental principles guided the data model:
\begin{enumerate}
  \item \textbf{Multi-tenant isolation via RLS}: every table carrying application
    data contains an \texttt{org\_id} column. RLS policies guarantee a query can
    only access the current organisation's rows (via \texttt{current\_org\_id()}).
  \item \textbf{Hypertables for time-series data}: \texttt{metrics} and
    \texttt{change\_events} are TimescaleDB hypertables partitioned automatically
    by time chunks. TimescaleDB applies columnar compression after 7 days,
    reducing storage by a factor of 10--20.
\end{enumerate}

\subsection{Key Tables by Domain}

\textbf{Authentication \& Management:} \texttt{organizations}, \texttt{users},
\texttt{servers}, \texttt{ca\_config}, \texttt{agent\_certs},
\texttt{enrollment\_tokens}, \texttt{refresh\_token\_blacklist}, \texttt{audit\_log}.

\textbf{Metrics \& Anomalies:} \texttt{metrics} (hypertable), \texttt{anomaly\_events},
\texttt{alert\_rules}, \texttt{notifications}.

\textbf{AI:} \texttt{ai\_explanations}, \texttt{server\_forecasts}, \texttt{rag\_documents},
\texttt{runbooks}.

\textbf{Change Correlation:} \texttt{change\_events} (hypertable),
\texttt{root\_cause\_analyses}.

\textbf{Splunk:} \texttt{splunk\_incidents}, \texttt{splunk\_ai\_explanations},
\texttt{ai\_feedback}, \texttt{org\_config}.

\section{Key Processing Flows}

\subsection{Change Correlation Flow}

\begin{enumerate}
  \item Correlation Engine receives \texttt{anomaly.detected.*} message;
  \item Queries \texttt{change\_events} in window $[t - 30\text{min}, t]$
    for the same server;
  \item Scores each change: $s = e^{-0.1 \cdot \Delta t} \times w_{\text{type}}$,
    where $\Delta t$ is minutes before anomaly and $w_{\text{type}} \in [0.3, 1.0]$;
  \item Top 3 changes sent to Ollama with anomaly context;
  \item Result persisted in \texttt{root\_cause\_analyses}, published to
    \texttt{rootcause.ready.\{org\_id\}}.
\end{enumerate}

\section{Security: STRIDE Threat Model}

\begin{table}[H]
\centering
\caption{STRIDE analysis and countermeasures in \caimp{}}
\label{tab:stride}
\small
\begin{tabularx}{\textwidth}{lXX}
\toprule
\textbf{Threat} & \textbf{Example} & \textbf{Countermeasure} \\
\midrule
Spoofing       & Malicious agent sending forged metrics     & mTLS ECDSA P-256; JWT \\
Tampering      & Metric modification in transit             & TLS 1.3; AEAD integrity \\
Repudiation    & Deny an administrative action              & Immutable audit\_log \\
Info Disclosure& Access to another org's data               & RLS PostgreSQL; Docker bridge \\
DoS            & Metric flood from compromised agent        & Nginx rate limiting \\
EoP            & Viewer accessing admin routes              & RBAC role check per endpoint \\
\bottomrule
\end{tabularx}
\end{table}

\section*{Conclusion}

This chapter presented the architectural foundations of \caimp{}: the microservices
design, four-layer organisation, NATS JetStream event bus, relational data model
and key processing flows. This architecture ensures the separation of concerns
necessary for platform scalability and maintainability. The next chapter justifies
the technology choices that give this architecture substance.
